\section{MOTIVATION}

Recent LLM-based unit test generation approaches, such as SymPrompt~\cite{ryan2024}, HITS~\cite{wang2024c}, Panta~\cite{gu2025a}, and JunitGenie~\cite{liao2025}, have achieved remarkable progress in automating test generation for complex methods, significantly reducing the manual effort required from developers. However, when the reachability of a target branch depends on \textit{object states established by prior method calls}, existing approaches still struggle to guide LLMs in constructing the appropriate external context, even when augmented with CFG-based path analysis.

The fundamental limitation is that these approaches typically supply the \textit{entire method body} as the LLM's input context. In real-world projects, however, the conditions necessary to reach a target branch are often not visible within the method body itself—they are implicit in a sequence of initialization and state-setting operations that must occur \textit{before} the method is invoked~\cite{baldoni2018surveysymbolicexecutiontechniques}. The irrelevant statements in the method body not only consume the context window but also obscure these cross-method dependencies, leaving the LLM unable to infer the required preconditions.

\textbf{A Motivating Example.}
Consider the \texttt{addValue()} method from Apache Commons 
Math~\cite{commonsMath}, shown in Listing~\ref{lst:motivation}.

\begin{lstlisting}[language=Java, 
                   caption={Target branch example from Apache Commons Math},
                   label={lst:motivation}]
// org.apache.commons.math3.stat.descriptive.SummaryStatistics
public void addValue(double value) {
    sumImpl.increment(value);
    ...
    
    if (meanImpl != mean) {   // target branch
        meanImpl.increment(value);
    }
    ...
    n++;
}
\end{lstlisting}

Suppose a test generation tool identifies that the branch at line~9 
(\texttt{meanImpl != mean}) has not been covered and attempts to generate 
a test to reach it. Existing approaches typically provide the entire method 
body alongside CFG path information to the LLM. However, the LLM faces a 
critical inferential barrier: the condition \texttt{meanImpl != mean} can 
only be satisfied if \texttt{setMeanImpl()} was called \textit{prior} to 
\texttt{addValue()}—a dependency that is entirely absent from the method 
body. Consequently, the LLM tends to generate a test such as:

\begin{lstlisting}[language=Java, caption={Typical LLM-generated test (target branch unreachable)}]
// meanImpl == mean always holds; target branch is never triggered
SummaryStatistics stats = new SummaryStatistics();
stats.addValue(1.0);
\end{lstlisting}

This test compiles and executes without error, yet the target branch is 
never triggered. The correct test requires establishing the necessary 
precondition \textit{before} invoking \texttt{addValue()}:

\begin{lstlisting}[language=Java, caption={Correct test (target branch reachable)}]
// Establish precondition via setMeanImpl() before calling addValue()
SummaryStatistics stats = new SummaryStatistics();
stats.setMeanImpl(new CustomMeanImpl()); // makes meanImpl != mean
stats.addValue(1.0);                     // target branch now reachable
\end{lstlisting}

This example illustrates two core challenges that existing approaches fail 
to address: \textbf{(1)} the method body contains numerous statements 
irrelevant to the target branch, which dilute the LLM's attention and 
impede precise reasoning; and \textbf{(2)} the true precondition for 
reaching the target branch is a cross-method dependency invisible within 
the method body itself.

\textbf{Our Approach.}
To address these challenges, we propose two complementary mechanisms that 
reduce LLM reasoning complexity from the perspectives of 
\textit{information reduction} and \textit{condition explicitness}:

\begin{itemize}[leftmargin=*]
    \item \textbf{Backward Slicing.} Starting from the target branch condition
    \texttt{meanImpl != mean}, we perform backward slicing to trace all
    statements with data or control dependencies on \texttt{meanImpl} and
    \texttt{mean}, including field declarations, default assignments in the
    constructor, and the definition of \texttt{setMeanImpl()}. This eliminates
    six irrelevant \texttt{increment} calls and other unrelated branches from
    the context, allowing the LLM to focus on statements that govern branch
    reachability.

    \item \textbf{Constraint Hints.} Building on the slice, we explicitly
    provide the LLM with the path condition and how it can be established,
    e.g., \emph{``To trigger the target branch, \texttt{meanImpl != mean} must
    hold; this can be achieved by calling \texttt{setMeanImpl(customImpl)}
    before \texttt{addValue()}.''} Rather than requiring the LLM to infer
    implicit cross-method dependencies from source code, this mechanism
    directly guides the LLM toward constructing the correct initialization
    sequence.
\end{itemize}

Together, these two mechanisms systematically reduce the difficulty of 
reasoning about complex branch conditions, enabling LLMs to generate tests 
that reach previously uncovered paths.